% !PRINTED 501
\weberpage{501}
\begin{center}
\S. 154.\\[4pt]
\gesperrt{Divisoren der Gruppen. Nebengruppen und conjugirte Gruppen.}
\end{center}

Wir haben im \S.\,152 gesehen, dass eine Function von $m$ unabh\"angigen Ver\"anderlichen immer eine Permutationsgruppe von $m$ Ziffern bestimmt, dass aber dies nicht immer zutrifft, wenn an Stelle der unabh\"angigen Ver\"anderlichen bestimmte Gr\"ossen gesetzt werden. Durch die \gesperrt{Galois}'sche Theorie wird aber dieser Unterschied ausgeglichen oder wenigstens auf seinen Kern zur\"uckgef\"uhrt. Es sei jetzt $F(x) = 0$ irgend eine Gleichung $m^{\text{ten}}$ Grades in einem gegebenen K\"orper $\Omega$ mit den von einander verschiedenen Wurzeln $\alpha, \alpha_1, \alpha_2 \ldots \alpha_{m-1}$.

Den Normalk\"orper
\[
\Omega(\varrho) = \Omega(\alpha, \alpha_1, \alpha_2 \ldots \alpha_{m-1})
\]
wollen wir mit $N$ bezeichnen.

Unter der \gesperrt{Galois}'schen Gruppe $P$ des K\"orpers $N$ oder der Gleichung $F(x) = 0$ k\"onnen wir (nach \S.\,146, 149) entweder die Gruppe der Substitutionen des K\"orpers $N$ oder auch die damit isomorphe Permutationsgruppe der Indices der $\alpha$ verstehen.

Es sei $p$ der Grad von $P$ und wir wollen die Operationen von $P$ (seien es nun Substitutionen von $N$ oder Permutationen der $\alpha$) mit
\weq{1}{P = \pi,\ \pi_1,\ \pi_2, \ldots \pi_{p-1}}
bezeichnen. Wenn das System $Q$ der in $P$ enthaltenen Permutationen
\weq{2}{\varkappa,\ \varkappa_1,\ \varkappa_2, \ldots \varkappa_{q-1}.}
% Source check: the faint terminal mark in (2) is the printed period.
f\"ur sich eine Gruppe ausmacht, wenn also je zwei der Elemente von $Q$, mit einander componirt, wieder ein Element aus $Q$ ergeben, so heisst die Gruppe $Q$ ein \gesperrt{Theiler} oder \gesperrt{Divisor von} $P$\footnote{Auch \gesperrt{Untergruppe} genannt.}).

Wir leiten zun\"achst einen wichtigen allgemeinen Satz \"uber die Theiler einer Gruppe her.

Wenn der Grad $q$ des Theilers kleiner ist als der Grad $p$ von $P$, so nehmen wir ein nicht in $Q$ enthaltenes Element $\pi_1$ in $P$ und bilden das System
\weq{3}{Q\,\pi_1 = \varkappa\,\pi_1,\ \varkappa_1\,\pi_1 \ldots \varkappa_{q-1}\,\pi_1,}
%CONT>>
